% !PRINTED 113
\weberpage{113}

Accordingly, if one applies the formulae § 31, (6), (11) in (2), one
finds for the linear fundamental transformations of
$\eta(\omega)$
\weq{4}{\eta(\omega+1) = e^{\frac{\pi i}{12}}\eta(\omega),}
\weq{5}{\eta\left(-\frac{1}{\omega}\right) = \sqrt{-i\omega}\,\eta(\omega),}
of which the first also follows immediately from the series representation
(3), while the second cannot be proved so easily by direct
transformation of the series.

We pass on to the consideration of the two principal transformations
of the second order of the $\eta$-function.

From § 32, (11) one obtains, by differentiation and setting
$u$ equal to zero,
\[ 2\vartheta_{01}(0, 2\omega)\eta(2\omega)^3 = \vartheta_{10}\eta(\omega)^3, \]
hence, if one raises this to the square and applies § 32, (8):
\[ 4\vartheta_{00}\vartheta_{10}\vartheta_{01}\eta(2\omega)^6 = \vartheta_{10}^3\eta(\omega)^6, \]
and, with the application of (2) and extraction of the cube root:
\weq{6}{2\eta(2\omega)^2 = \vartheta_{10}\eta(\omega).}

If in (6) one replaces $\omega$ by $\dfrac{\omega}{2}$, it follows that
\[ 2\eta(\omega)^2 = \vartheta_{10}\left(0, \frac{\omega}{2}\right)\eta\left(\frac{\omega}{2}\right), \]
and if one squares and applies the formula § 32 (9):
\[ 2\eta(\omega)^4 = \eta\left(\frac{\omega}{2}\right)^2\vartheta_{00}\vartheta_{10}. \]

If one multiplies both sides by $\vartheta_{01}$, it follows
by (2):
\weq{7}{\eta\left(\frac{\omega}{2}\right)^2 = \vartheta_{01}\eta(\omega)}
and by changing $\omega$ into $\omega+1$ [(4) and § 31,
(6)]:
\weq{8}{e^{-\frac{\pi i}{12}}\eta\left(\frac{\omega+1}{2}\right)^2 = \vartheta_{00}\eta(\omega).}
